% Diagrama del proceso de investigación estructurado en dos filas continuas
% Fila 1: Etapas 1-4 (OE1: Analizar)
% Fila 2: Etapas 5-7 (OE2: Comparar, OE3: Evaluar, OE4: Evaluar)
\begin{tikzpicture}[
    font=\footnotesize,
    etapaFila1/.style={
        rectangle,
        rounded corners=4pt,
        draw=ugrAzul,
        fill=white,
        line width=0.8pt,
        text width=3.30cm,
        minimum height=2.65cm,
        align=center,
        inner sep=4pt
    },
    etapaFila2/.style={
        rectangle,
        rounded corners=4pt,
        draw=ugrAzul,
        fill=white,
        line width=0.8pt,
        text width=4.55cm,
        minimum height=2.65cm,
        align=center,
        inner sep=4pt
    },
    flechaEtapa/.style={
        -{Stealth[length=2.5mm, width=1.8mm]},
        line width=1.0pt,
        color=ugrAzul!85
    }
]

    % ================= FILA 1: OE1 (ETAPAS 1 A 4) =================
    % 1. Adquisición (OE1)
    \node[etapaFila1] (e1) at (0, 0) {%
        {\scriptsize\bfseries\color{ugrAzul} 1. Adquisición}\\[2pt]
        {\tiny\bfseries\color{ugrVerde} [OE1: Analizar]}\\[3pt]
        {\scriptsize Descarga de Twitter Coello y colecciones Reddit}\\[3pt]
        \includegraphics[height=0.45cm]{logos/huggingface.png}\\[2pt]
        {\tiny\texttt{huggingface\_hub} $\cdot$ \texttt{requests}}
    };

    % 2. Traducción (OE1)
    \node[etapaFila1] (e2) at (3.85, 0) {%
        {\scriptsize\bfseries\color{ugrAzul} 2. Traducción}\\[2pt]
        {\tiny\bfseries\color{ugrVerde} [OE1: Analizar]}\\[3pt]
        {\scriptsize MarianMT EN$\to$ES local con decodificación voraz}\\[3pt]
        \includegraphics[height=0.45cm]{logos/pytorch.png}\\[2pt]
        {\tiny MarianMT (\texttt{opus-mt-en-es})}
    };

    % 3. EDA (OE1)
    \node[etapaFila1] (e3) at (7.70, 0) {%
        {\scriptsize\bfseries\color{ugrAzul} 3. EDA}\\[2pt]
        {\tiny\bfseries\color{ugrVerde} [OE1: Analizar]}\\[3pt]
        {\scriptsize Clases, longitudes, vocabulario y marcadores}\\[3pt]
        \includegraphics[height=0.45cm]{logos/pandas.png}\\[2pt]
        {\tiny \texttt{pandas} $\cdot$ \texttt{matplotlib}}
    };

    % 4. Preprocesamiento (OE1)
    \node[etapaFila1] (e4) at (11.55, 0) {%
        {\scriptsize\bfseries\color{ugrAzul} 4. Preproceso}\\[2pt]
        {\tiny\bfseries\color{ugrVerde} [OE1: Analizar]}\\[3pt]
        {\scriptsize Anonimización determinista y partición 70/10/20}\\[3pt]
        \includegraphics[height=0.45cm]{logos/spacy.png}\\[2pt]
        {\tiny \texttt{spaCy} $\cdot$ \texttt{Python re}}
    };

    % ================= FILA 2: OE2, OE3, OE4 (ETAPAS 5 A 7) =================
    % 5. Baselines (OE2)
    \node[etapaFila2, draw=ugrNaranja, fill=ugrNaranja!4] (e5) at (0.65, -3.45) {%
        {\scriptsize\bfseries\color{ugrAzul} 5. Baselines Léxicos}\\[2pt]
        {\tiny\bfseries\color{ugrNaranja!80!black} [OE2: Comparar]}\\[3pt]
        {\scriptsize TF-IDF con Regresión Logística y SVM lineal}\\[3pt]
        \includegraphics[height=0.45cm]{logos/scikitlearn.png}\\[2pt]
        {\tiny \texttt{scikit-learn} $\cdot$ Submuestreo / balanceo}
    };

    % 6. Ajuste Fino BETO (OE3)
    \node[etapaFila2, draw=ugrAzul, fill=filaResaltada] (e6) at (5.80, -3.45) {%
        {\scriptsize\bfseries\color{ugrAzul} 6. BETO Contextual}\\[2pt]
        {\tiny\bfseries\color{ugrAzul} [OE3: Evaluar]}\\[3pt]
        {\scriptsize Fine-tuning Transformer (BETO, 110M params, FP16)}\\[3pt]
        \includegraphics[height=0.45cm]{logos/pytorch.png}\\[2pt]
        {\tiny \texttt{PyTorch} $\cdot$ \texttt{transformers} (AdamW)}
    };

    % 7. Evaluación (OE4)
    \node[etapaFila2, draw=ugrVerde, fill=ugrVerde!6] (e7) at (10.95, -3.45) {%
        {\scriptsize\bfseries\color{ugrAzul} 7. Evaluación y Auditoría}\\[2pt]
        {\tiny\bfseries\color{ugrVerde} [OE4: Evaluar]}\\[3pt]
        {\scriptsize $F_1$-macro, AUC, bootstrap pareado y robustez}\\[3pt]
        \includegraphics[height=0.45cm]{logos/scipy.png}\\[2pt]
        {\tiny \texttt{scipy} $\cdot$ \texttt{scikit-learn}}
    };

    % Flechas horizontales Fila 1
    \draw[flechaEtapa] (e1) -- (e2);
    \draw[flechaEtapa] (e2) -- (e3);
    \draw[flechaEtapa] (e3) -- (e4);

    % Conexión entre Fila 1 y Fila 2: desde etapa 4 hacia etapa 5
    \draw[flechaEtapa] (e4.south) -- ++(0, -0.35) -| (e5.north);

    % Flechas horizontales Fila 2
    \draw[flechaEtapa] (e5) -- (e6);
    \draw[flechaEtapa] (e6) -- (e7);

\end{tikzpicture}
